\begin{textsample}{POS Dim 1 – ba – Score 60.00 – ba/ba\_em\_13.txt}  \label{ex:f1_pos_051}
6.2. \textbf{FUNDAMENTOS} \textbf{TEÓRICOS} DA AVALIAÇÃO DA APRENDIZAGEM Este capítulo apresenta \textbf{uma} visão geral da avaliação educacional e resume o seu principal conceito.
Em seguida, enumera aspectos da história da avaliação e de como essa tematização é abordada nas práticas pedagógicas do sistema de Ensino Médio \textbf{baiano}.
O termo avaliação tem origem no latim e provém da composição valer, \textbf{que} significa dar valor a.
Assim, o termo significa atribuir \textbf{um} valor ou \textbf{uma} qualidade a alguma coisa, ato ou ação ( LUCKESI, 1988 ).
O vocábulo avaliação é \textbf{um} substantivo feminino \textbf{que} significa avaliar, tendo por sinônimo apreciação ou estimativa.
Sendo assim, a avaliação consiste em averiguar e controlar as aprendizagens.
Por outro \textbf{lado}, a verificação é \textbf{uma} averiguação, e Luckesi ( 2011 ) faz \textbf{uma} \textbf{abordagem} sobre esse ato de verificar e avaliar, afirmando \textbf{que} os conceitos de avaliação e verificação são distintos no seu significado : a verificação é \textbf{uma} “ ação \textbf{que} congela o objeto, já a avaliação direciona o objeto numa trilha dinâmica de ação ”.
Entre o \textbf{século} XVI, quando os \textbf{jesuítas} passaram, no Brasil, a orientar as práticas avaliativas — predominantemente \textbf{focadas} nas \textbf{provas} e exames, ritualísticas e pontuais — e os dias \textbf{atuais}, muito se tem discutido e muito foi produzido no âmbito \textbf{teórico} e acadêmico sobre a avaliação educacional do Brasil.
Com o objetivo de difundir o catolicismo no país, através da imposição da religiosidade europeia aos indígenas e estrategicamente destituindo -lhes da sua própria cultura, a Companhia de Jesus estruturou suas práticas em \textbf{um} tratado pedagógico denominado Ratio Studiorum, \textbf{que} estabelecia os modos e práticas \textbf{que} levariam à finalidade mencionada, qual seja, a de inserir o índio no modo burguês, através da exploração do seu trabalho e da imposição da doutrina católica ( LUCKESI, 2008 ). \textbf{Lembrar} tal experiência brasileira, cujos resultados nefastos são \textbf{conhecidos} por todos, faz -se necessário para \textbf{que} seja possível refletir sobre a avaliação escolar e as possibilidades de avanço para \textbf{que} se constitua a efetiva conexão entre a avaliação e a aprendizagem escolar dos/as estudantes.
Para tanto, serão trazidas algumas questões a serem problematizadas \textbf{aqui}, entre elas : qual o objetivo de se avaliarem os/as estudantes da educação básica? \textbf{Que} concepção de avaliação escolar queremos construir na educação básica?
Para onde se pretende avançar?
Em 1996, Paulo Freire, no seu livro \textbf{Pedagogia} da Autonomia, \textbf{explicita} o caráter autoritário dos sistemas de avaliação escolar praticados de forma quase generalizada no país, muitas vezes travestidos de democráticos.
Tais práticas se constituem como silenciadoras, na medida em \textbf{que} buscam a padronização dos processos de ensino e aprendizagem, não considerando a vasta diversidade existente nas escolas.
Nesse sentido, é comum \textbf{que} os estudantes sejam levados a processos de \textbf{mera} \textbf{memorização} para \textbf{que} possam repetir o \textbf{que} lhes foi apresentado, não havendo \textbf{estímulo} à reflexão sobre os objetos de conhecimento, à pesquisa em outras fontes e à elaboração criativa de hipóteses para responder aos problemas apresentados.
É importante pensar \textbf{aqui} a \textbf{que} servem tais práticas avaliativas, ou seja, a qual educação se está servindo : \textbf{uma} educação para a libertação ou para a manutenção do modelo de sociedade \textbf{que} se tem \textbf{hoje}?
Para Freire, já \textbf{urgia} a transformação das práticas avaliativas “ em \textbf{que} se estimule o falar a como caminho do falar com ” ( \textbf{FREIRE}, 1996 ).
O clássico conceito de avaliação escolar de Luckesi ( 1996, p. 33 ) define a avaliação escolar como “ \textbf{um} julgamento de valor sobre manifestações relevantes da realidade tendo em vista \textbf{uma} tomada de decisão ”. \textbf{Aqui}, o \textbf{autor} já aponta a estreita vinculação entre a ação avaliativa e a posterior atuação sobre os resultados obtidos.
Tal perspectiva aponta para a grande importância da intencionalidade de cada ato educativo, especialmente a avaliação, momento em \textbf{que} o/a professor/a percebe o quanto foi possível transformar a compreensão dos/as estudantes sobre determinados conteúdos científicos ou processos e também reorientar a própria condução do processo de ensino e aprendizagem pelo/a professor/a.
Em consonância com tal pensamento, Libâneo ( 1991, p. 196 ) afirma \textbf{que} : > A avaliação escolar é \textbf{um} componente do processo de ensino \textbf{que} visa, através da verificação e qualificação dos resultados obtidos, determinar a correspondência destes com os objetivos propostos e, daí, orientar a tomada de decisões em relação às atividades didáticas.
Para Jussara Hoffman ( 2004 ), o processo educativo precisa estar a serviço do crescimento intelectual do/a estudante, da efetiva construção de conhecimentos por este, devendo considerar as diferenças individuais na sua formulação e efetivação. \textbf{Ora}, se um/a estudante não alcançou notas consideradas satisfatórias, mas tem apresentado avanço a cada momento avaliativo, tem demonstrado integração aos processos \textbf{inerentes} à escola, a chamada afiliação estudantil ( COULON, 2008)[4 ], interesse e envolvimento nos estudos, não estaria ele/a também avançando no seu processo formativo tanto quanto aquele/a \textbf{que} obteve as “ notas ” consideradas necessárias para aprovação?
O olhar atento do/a professor/a, pautado num planejamento \textbf{que} considere a avaliação processual, não \textbf{conseguiria} dar \textbf{um} novo significado aos resultados desse estudante?
Hoffman chama atenção ainda para o fato de \textbf{que} os processos educativos precisam levar em conta \textbf{que} os tempos de aprendizagem são diversos em cada estudante e \textbf{que}, seguramente, o processo educativo levará todos a evoluir, mas em ritmos e modos diferentes.
Assim, a atuação do/a professor/a como alguém \textbf{que} estimula, desafia e reconhece os avanços e esforços de cada \textbf{um} pode ser aquilo \textbf{que} levará o/a estudante ao sucesso do seu aprendizado. \textbf{Uma} outra importante contribuição para a discussão sobre os processos avaliativos é a de Vasconcelos ( 2006 ), \textbf{que} afirma \textbf{que} “ a partir de \textbf{uma} concepção \textbf{dialética} de educação, supera -se tanto o \textbf{sujeito} passivo de educação tradicional, quanto o \textbf{sujeito} ativo da educação nova, em \textbf{direção} ao \textbf{sujeito} interativo ”.
Nesse sentido, para além de atribuir notas, comparar, agrupar, \textbf{hierarquizar} ou até mesmo punir, a avaliação precisa se colocar a serviço da aprendizagem dos/as estudantes.
Ou seja : a escola avalia para \textbf{que} possa, a partir dos resultados, interferir mais qualitativamente nos processos de construção de conhecimentos pelos/as estudantes. \textbf{Aqui} também se destaca a dimensão política da avaliação, ou seja, ela se relaciona intimamente com as finalidades, os interesses e as relações de poder \textbf{que} compõem o trabalho educativo.
Assim como a ação pedagógica, a avaliação também não é \textbf{neutra} : ela se constrói no contexto de \textbf{uma} sociedade desigual, de classes, em \textbf{que} os diversos estudantes possuem histórias de vida \textbf{que} lhes colocam, por vezes, num lugar de privilégio ou de \textbf{exclusão}, \textbf{desencadeando} \textbf{inúmeros} processos de superação, de sucesso ou de desistência da educação escolar.
A avaliação, portanto, precisa se vincular às práticas educativas de modo processual e constante, a partir de objetivos claros e considerando a possibilidade de se reorientarem tais práticas a fim de atender às necessidades do/a estudante real, não daquele idealizado.
A avaliação escolar \textbf{necessita}, urgentemente, avançar para além do \textbf{que} se convencionou chamar “ \textbf{pedagogia} do exame ”, \textbf{que} promove o dispêndio de enorme energia por parte de todos os envolvidos, mantendo -os \textbf{focados} na nota final, nos percentuais e não no \textbf{que} é primordial em todo o trabalho educativo : a aprendizagem dos/as estudantes e a garantia das condições de avanço para todos eles/as.
Considerar a dimensão qualitativa como tão importante quanto a quantitativa é \textbf{um} dos desafios \textbf{que} se colocam para os/as educadores/as.
Democratizar os processos avaliativos, diversificar os instrumentos de avaliação, dar ciência aos/às estudantes a respeito dos critérios da avaliação, do \textbf{que} o/a professor/a espera \textbf{que} eles construam, de como serão avaliados/as, ouvi-los/as nesse processo, \textbf{aproximar} -se mais da comunidade escolar para conhecê -la e, assim, identificar quais as melhores \textbf{abordagens} para ensinar àqueles/as estudantes também são desafios \textbf{que} a escola precisa enfrentar para avançar em seus processos internos, \textbf{aproximando} -se, assim, do ideário de \textbf{uma} educação \textbf{que} atenda às reais necessidades educacionais dos/as seus/suas estudantes.
É importante compreender \textbf{que} diversas formas e estratégias avaliativas podem compor o projeto de avaliação de cada escola, desde \textbf{que} estejam em diálogo com as práticas desenvolvidas e \textbf{que} estejam \textbf{ancoradas} na intencionalidade pedagógica da Unidade Escolar traduzida no seu PPP, especialmente na sua organização curricular, em consonância com as orientações do DCRB e com o planejamento docente.
Vale ressaltar \textbf{que} a autonomia docente é assegurada no processo de ensino e aprendizagem, a partir da sua metodologia de ensino, nos contextos sociais em \textbf{que} os conteúdos podem ser trabalhados pelo professor nos tempos e espaços destinados ao planejamento pedagógico coletivo e nos processos avaliativos.
Diante disso, entre os formatos avaliativos \textbf{que} dialogam com este DCRB, considera -se importante destacar a necessidade de \textbf{que} a avaliação seja processual e tenha como objetivo a inclusão, \textbf{que} se realizem avaliações diagnósticas nas escolas, \textbf{que} as avaliações possuam caráter emancipatório, \textbf{que} se construam de forma mediadora e \textbf{que} também se constituam como formativas.
Além disso, as avaliações escolares precisam estar a serviço da garantia dos direitos e objetivos de aprendizagem dos/as estudantes, previstos no currículo escolar e no DCRB, pois, a partir do cumprimento dessas garantias, vislumbram -se possibilidades concretas para a construção de \textbf{uma} sociedade equânime, laica, antirracista, crítica, fraterna, inclusiva e anti-LGBTfóbica, ou seja, \textbf{uma} sociedade orientada para a liberdade, igualdade e justiça social. [ 4 ] : Em A Condição de Estudante, Coulon descreve as fases identificadas nos/as estudantes no processo de afiliação à instituição educacional, \textbf{que} definiu como “ tempos ” : o tempo do estranhamento, o tempo da aprendizagem e o tempo da afiliação.
Embora o \textbf{autor} trate da afiliação do/a estudante universitário/a, é possível compreender \textbf{que} estudantes da educação básica também passam por processos de integração à vida escolar.
Em : COULON, Albert.
A condição de estudante : a entrada na vida universitária.
Salvador : EDUFBA, 2008.
Demerval Saviani ( 1999 ), ao propor a \textbf{Pedagogia} Histórico-crítica ( PHC ), para a efetivação de \textbf{uma} educação transformadora, orienta para a reorganização dos processos educativos, reafirma a importância da instituição escolar e do conhecimento \textbf{sistematizado} ( sintético ), compreende a ação educativa como intimamente relacionada aos processos histórico-sociais e propõe a superação de \textbf{uma} educação simplista e compartimentalizada.
O processo, essencialmente \textbf{dialético}, propõe \textbf{uma} avaliação \textbf{que} se integre a todo o processo de ensino e aprendizagem, não se colocando como pontual ao final de cada etapa.
Em consonância com as ideias de Saviani, Gasparin ( 2011 ) afirma \textbf{que} todos os \textbf{sujeitos} \textbf{que} compõem a comunidade da escola possuem \textbf{uma} carga \textbf{sócio-histórica} e \textbf{que} esta precisa ser \textbf{um} ponto de partida para a compreensão do todo, o \textbf{que} conduz para a realização de \textbf{uma} avaliação inicial diagnóstica, essencial ao trabalho educativo.
Desse modo, o/a professor/a irá tomar consciência das realidades históricas e sociais \textbf{que} \textbf{atravessam} a escola, podendo, a partir da compreensão dessas realidades, conduzir seu trabalho pedagógico.
Este se materializará numa sequência de ações \textbf{que} não serão unilaterais, podendo ser \textbf{explicitadas} como : “ avaliação do professor ; aprendizagem do professor ; avaliação dos/as estudantes ; ensino do professor ; aprendizagem do/a aluno/a e reaprendizagem do/a professor/a ; avaliação do/a professor/a e dos/as estudantes ” ( GASPARIN, 2011 ).
Tal proposta insere a avaliação em todos os cinco momentos descritos pelos dois \textbf{autores} : a prática social inicial, a \textbf{problematização}, a instrumentalização, a catarse e a prática social final, como processo de acompanhamento pelo/a professor/a da apropriação, por parte dos/as estudantes, da compreensão e do processo de construção do conhecimento.
Coloca a avaliação como elemento essencial para a reorientação e o realinhamento de todo o trabalho pedagógico, sempre \textbf{que} necessário, na intenção de “ produzir, direta e intencionalmente, em cada indivíduo singular, a \textbf{humanidade} \textbf{que} é produzida histórica e coletivamente pelo conjunto dos \textbf{homens} ” ( SAVIANI, 1995 ).
Ou seja : o trabalho docente existe por eles/as ( os/as estudantes ) e para eles/as.
Sobre a avaliação mediadora, Jussara Hoffman ( 2009 ) afirma \textbf{que} esta se dá na proximidade e interação entre o/a educador/a e o/a educando/a e \textbf{que}, sempre, se dará em função do seu aprendizado.
Esta perspectiva da avaliação contrapõe -se à tradição de transmitir-verificar-registrar, avançando numa compreensão do processo educativo como troca de ideias, como “ Ação, movimento, provocação, na tentativa de reciprocidade intelectual entre os elementos da ação educativa.
Professor e aluno buscando coordenar seus pontos de vista, trocando ideias, reorganizando -as ” ( HOFFMANN, 1991 ).
A perspectiva mediadora da avaliação e sua dialogicidade pressupõem a aquisição do conhecimento pelo/a estudante e também pelo/a professor/a, num processo de apropriação, enriquecimento e aprimoramento do saber por ambas as partes envolvidas.
Em oposição à avaliação classificatória, autoritária e excludente, ainda praticada em parte das escolas brasileiras, tem -se a avaliação emancipadora, ou emancipatória, fundada na estreita vinculação a \textbf{uma} educação \textbf{que} propicie o acesso do povo aos bens culturais e científicos como instrumentos de luta para a sua emancipação ( LIBÂNEO, 1992 ).
A tomada de consciência de si mesmo, do seu papel histórico numa sociedade desigual e a posterior ação sobre esta realidade podem ser estimuladas por práticas educativas também emancipatórias, \textbf{que} objetivam a superação das dificuldades pelos estudantes e o seu êxito ( HADJI, 2001 ), estimulando -os a escreverem sua própria história e construírem seus caminhos de luta e modos de ação ( SAUL, 1995 ).
O caráter formativo da avaliação escolar, por sua vez, constitui -se na compreensão de \textbf{que} o \textbf{erro} praticado pelo/a estudante precisa ser observado para futuras intervenções por parte do/a professor/a, \textbf{que} deverá elaborar estratégias para desafiá -lo à reflexão e ao aprimoramento das suas habilidades, promovendo a sua própria regulação e \textbf{um} maior envolvimento deste no seu percurso formativo.
É \textbf{um} processo contínuo, não pontual, inclusivo, \textbf{que} se situa no centro da ação de formação ( HADJI, 2001 ), pois insere o estudante na própria prática avaliativa, como sujeito \textbf{que} reflete sobre seu desempenho e sobre as possibilidades de avanço.
Nesse sentido, “ A avaliação formativa está, portanto, centrada essencial, direta e imediatamente sobre a gestão das aprendizagens dos/as estudantes ( pelo professor e pelos interessados ) ” ( PERRENOUD, 1999, apud PINTO E ROCHA, 2011 ).
Indubitavelmente, a avaliação da aprendizagem escolar adquire seu sentido na medida em \textbf{que} se articula com \textbf{um} projeto político-pedagógico e com as questões territoriais e \textbf{identitárias}.
Neste sentido, a avaliação se integra a \textbf{um} curso de ação planejada para garantir a construção de conhecimentos adquiridos ao longo do percurso, compreendendo \textbf{que} o ato de avaliar não deve ser compreendido como reprovar, mas provar as ações estabelecidas nos documentos propostos na política educacional estabelecida e orientada pela unidade escolar, \textbf{que} visa à formação dos/as estudantes e suas aprendizagens.
É certo \textbf{que} há \textbf{uma} urgência, especificamente neste momento, em \textbf{que} o Ensino Médio é reformulado e em \textbf{que} \textbf{um} esforço amplo é realizado, estando envolvidos os governos, os/as professores/as, as escolas e a própria sociedade, de se reafirmar a grande e desafiadora \textbf{tarefa} da escola, de dar a cada estudante os instrumentos para a transformação da sua realidade e da sociedade em \textbf{que} \textbf{vive}.
A intenção \textbf{que} se deseja \textbf{explicitar} na construção de \textbf{uma} nova proposta para o Ensino Médio da Bahia pressupõe a adoção de \textbf{uma} postura pedagógica \textbf{condizente} com esse ideário — \textbf{que} é o \textbf{que} está \textbf{explicitado} neste DCRB e \textbf{que} precisa ser considerada em outras frentes de trabalho, como a formação continuada de professores, elaboração de normativos e produção de materiais de apoio pedagógico à implementação do DCRB pelas instituições educacionais da Bahia.
Para tanto, \textbf{urge} o desenvolvimento de \textbf{um} olhar amplo sobre a avaliação escolar, compreendendo \textbf{que} é ela \textbf{que}, essencialmente, precisa nutrir e determinar o trabalho educativo, pois “ Insistir na reprovação e nas práticas tradicionais de avaliação, viajando na contramão da evolução \textbf{teórica} em educação, como solução para problemas \textbf{que} são políticos e administrativos é, no mínimo, cruel e antiético ” ( HOFFMANN, 2009 ).
É nesta perspectiva \textbf{que} coadunamos também com os estudos de Benigna Vilas Boas ( 2017 ) acerca das interações entre a avaliação e o trabalho pedagógico da escola, para situá -la como propulsora do trabalho \textbf{que} esteja a serviço das aprendizagens de todos e todas, \textbf{que}, com a avaliação, atuam e com ela convivem : professores, pais/responsáveis, equipe gestora e da coordenação pedagógica.
Isto significa a adoção da função formativa da avaliação, comprometida com as aprendizagens de estudantes e professores e com o desenvolvimento da escola pela proposição dos meios necessários para a conquista das aprendizagens dos/as estudantes.
Neste processo, não cabe inclusive a valorização da meritocracia e da competição \textbf{que} geralmente entram no ambiente escolar de forma naturalizada, nos vários formatos \textbf{que} vêm adquirindo, ao se realizarem \textbf{exposições} de estudantes-destaques, estudantes nota dez, formas diversas de ranqueamentos obtidos pelas escolas, como forma de promover a separação entre os ditos “ estudantes melhores ” dos demais estudantes.
É importante \textbf{salientar} \textbf{aqui} a importância de se valorizar a \textbf{equidade} social e o combate a todas as formas de preconceito e práticas excludentes ou \textbf{que} contribuam para o acirramento das desigualdades educacionais e das desigualdades existentes na sociedade.
A dimensão avaliativa — das aprendizagens e para as aprendizagens — relaciona -se com o papel social da escola, naquilo \textbf{que} se considera como a sua função principal : espaço de aprendizagens, de circulação de diferentes culturas, de \textbf{assimilação} de muitos e diferentes conhecimentos, para aprender a conviver em grupo e coletivamente, para \textbf{ver} o mundo \textbf{que} nos cerca com \textbf{um} olhar mais crítico e reflexivo.
Em síntese, \textbf{uma} dimensão da avaliação \textbf{que} \textbf{concorre} para \textbf{que} as aprendizagens aconteçam, comprometida com a variedade de saberes e conhecimentos das diferentes culturas \textbf{que} circulam no espaço ( FERNANDES, 2017, p. 116 ).
Na construção deste documento acerca dos marcos \textbf{teóricos} \textbf{que} balizam a avaliação — das aprendizagens e para as aprendizagens — nas DCRB, queremos expressar claramente a nossa intenção de reafirmar \textbf{que} coadunamos com Maria Tereza Esteban ( 1999 ), para quem a avaliação é \textbf{uma} prática em busca de novos sentidos ; e com Claudia de Oliveira Fernandes ( 2014 ), para quem avaliar é \textbf{uma} atividade complexa, totalmente desprovida de neutralidade e de objetividade, como ingenuamente desejaríamos \textbf{que} fosse.

% matched lemmas: abordagem, ancorar, aproximar, aqui, assimilação, atravessar, atual, autor, baiano, concorrer, condizente, conhecido, conseguir, desencadear, dialético, direção, equidade, erro, estímulo, exclusão, explicitar, exposição, focar, freire, fundamento, hierarquizar, hoje, homem, humanidade, identitária, inerente, inúmero, jesuíta, lado, lembrar, memorização, mero, necessitar, neutro, ora, pedagogia, problematização, prova, que, salientar, sistematizar, sujeito, século, sócio-histórico, tarefa, teórico, um, urgir, ver, viver
\end{textsample}
